%% kapa: exam=1 round=29 year=2018 period=2 form=A subject=appraisal_law subject_ko=감정평가관계법규 seq=32 number=32
%% source: Q-Net 감정평가사 기출문제 (cst00309) artlSeq=5202808 file=extracted/1008/5202808/2018년 제29회 감정평가사 1차시험 기출문제/2018년 제29회 감정평가사 1차시험 기출문제/2교시 A형.pdf
%% conversion: kordoc 4.24.51 -> tools/gen1.py (auto); review=proofread-sonnet; flags=none
\begin{question}{32}
부동산등기법상 등기신청인에 관한 설명으로 옳지 않은 것은?
\begin{choices}
  \item 소유권보존등기 또는 소유권보존등기의 말소등기는 등기명의인으로 될 자 또는 등기명의인이 단독으로 신청한다.
  \item 상속, 법인의 합병, 그 밖에 대법원규칙으로 정하는 포괄승계에 따른 등기는 등기권리자가 단독으로 신청한다.
  \item 판결에 의한 등기는 승소한 등기권리자 또는 등기의무자가 단독으로 신청한다.
  \item 신탁재산에 속하는 부동산의 신탁등기는 위탁자가 단독으로 신청한다.
  \item 등기명의인표시의 변경이나 경정의 등기는 해당 권리의 등기명의인이 단독으로 신청한다.
\end{choices}
\end{question}
\answer{4}
